% Generated draft: requires mathematical and source review.
\documentclass[11pt]{article}
\usepackage[margin=25mm]{geometry}
\usepackage{fontspec}
\setmainfont{Latin Modern Roman}
\usepackage{amsmath,amssymb,mathtools,unicode-math}
\setmathfont{Latin Modern Math}
\usepackage{enumitem,needspace,xurl,hyperref}
\setlength{\parindent}{0pt}
\setlength{\parskip}{5pt}
\setlength{\emergencystretch}{3em}
\newcommand{\N}{\mathbb{N}}
\newcommand{\Z}{\mathbb{Z}}
\newcommand{\Q}{\mathbb{Q}}
\newcommand{\R}{\mathbb{R}}
\newcommand{\C}{\mathbb{C}}
\newcommand{\F}{\mathbb{F}}
\newcommand{\A}{\mathbb{A}}
\newcommand{\PP}{\mathbb{P}}
\newcommand{\E}{\mathbb{E}}
\newcommand{\eps}{\varepsilon}
\newcommand{\dd}{\,\mathrm{d}}
\begin{document}
\section*{UnsolvedMath ID: 30001171}
\textbf{Source problem code:} catalogue entry
\subsection{Definitions and mathematical statement}
The supplied record collects several questions rather than one fully specified conjecture. The formulations below retain their question form and explicitly identify missing hypotheses.

For a self-adjoint operator with discrete lowest eigenvalues, write \(\lambda_1\leq\lambda_2\) for its first two eigenvalues, counted with multiplicity, and call \(\lambda_2-\lambda_1\) its fundamental gap. A ground state is an eigenfunction associated with \(\lambda_1\). The operator domain and boundary conditions are part of the spectral problem and must be specified before these quantities are determined.

\begin{enumerate}
\item \textbf{Magnetic gap question.} Can a valid magnetic fundamental-gap inequality be obtained by subtracting a correction depending on the magnetic potential \(A\)? Schematically, the requested form is
\[
\lambda_2(A)-\lambda_1(A)\geq B-C(A),
\]
where \(B\) denotes the intended reference lower bound and \(C(A)\) the proposed correction. This is only a template: the supplied record does not specify the magnetic operator, its boundary conditions, the reference bound \(B\), or the permitted dependence and normalization of \(C(A)\). It therefore does not determine a quantified inequality.

\item \textbf{Ground-state log-concavity question.} Is the ground state of \(-\Delta\) log-concave? For a positive function \(u_1\) on a convex set \(\Omega\subset\R^n\), log-concavity means
\[
\log u_1((1-t)x+ty)\geq(1-t)\log u_1(x)+t\log u_1(y)
\quad(x,y\in\Omega,\ 0\leq t\leq1).
\]
This defines the property being asked about; it does not supply the missing class of domains or boundary conditions. The record does not explicitly state these hypotheses or specify a positive ground-state representative.

\item \textbf{Dirichlet eigenfunction-ratio question.} On a convex domain \(\Omega\), let \(u_1,u_2\) denote first and second eigenfunctions of the intended Dirichlet spectral problem. Wherever \(u_1>0\), put
\[
q(x)=\frac{u_2(x)}{u_1(x)}.
\]
Does this ratio attain its minimum and maximum on \(\partial\Omega\), as an analogue of the Hot Spots question? Since both Dirichlet eigenfunctions vanish on the boundary in the classical formulation, the quotient is not defined there by direct division. If an appropriate continuous extension \(\overline q\) to a compact closure exists and is the intended interpretation, the proposed assertion would be
\[
\exists x_-,x_+\in\partial\Omega:\qquad
\overline q(x_-)=\min_{\overline\Omega}\overline q,
\qquad
\overline q(x_+)=\max_{\overline\Omega}\overline q.
\]
The extension interpretation, its existence, the underlying differential operator, the dimension, and sufficient domain regularity are unresolved in the supplied formulation. The choice or quantification of \(u_2\) when \(\lambda_2\) is multiple is also unspecified.

\item \textbf{Robin gap question.} Using \(H_V=-\Delta+V\) as notation for the nonmagnetic Schr\"odinger operator suggested by the record, consider the Robin eigenvalue problem
\[
H_Vu=\lambda u\quad\text{in }\Omega,
\qquad
\frac{\partial u}{\partial\nu}=-\alpha u
\quad\text{on }\partial\Omega,
\]
where \(\nu\) is the outward unit normal. For a domain of finite diameter
\[
D=\operatorname{diam}(\Omega)
=\sup\{|x-y|:x,y\in\Omega\},
\]
is the fundamental gap minimized by taking \(V=0\) and letting the domain degenerate to a segment of diameter \(D\)? Can Lavine's methods be adapted to the one-dimensional Robin problem? The admissible domains and potentials, the range and constancy of \(\alpha\), the quantities held fixed, and the meaning of degeneration are not specified. In particular, the record does not identify the limiting interval boundary parameter or state a convergence assertion. Thus neither a precise variational infimum nor an interval comparison inequality is determined by the supplied text. The reference to Lavine's methods is also not accompanied by a specific result or citation.
\end{enumerate}
\subsection{Short English statement}
This record asks about magnetic corrections to spectral-gap bounds, log-concavity of Laplace ground states, boundary extrema of a Dirichlet eigenfunction ratio on convex domains, and minimization of Robin gaps by a zero-potential segment limit. Several hypotheses and boundary interpretations are missing. The supplied status is partially solved; this draft preserves that dated assessment without independently verifying it.
\subsection{Sources}
\begin{itemize}
\item Supplied local UnsolvedMath registry record, permanent ID \(30001171\), source problem code \texttt{catalogue entry}, title ``Magnetic, Robin, and Hot-Spots Spectral Gap Questions.'' The record belongs to the Oberwolfach Reports Open Problems collection and gives the proposed year as 2009. No individual proposer is identified.
\item Supplied source citation: \emph{Low Eigenvalues of Laplace and Schr\"odinger Operators} (2009), Oberwolfach Reports, pp.~355--428, DOI \texttt{10.4171/owr/2009/06}. \url{https://doi.org/10.4171/owr/2009/06}. The background describes these reports as documenting mathematical workshops and their open-problem discussions; no item-specific page is supplied.
\item Supplied literature source labeled ``fundamental gap theorem'': \url{https://doi.org/10.1090/S0894-0347-2011-00699-1}.
\item Supplied literature source labeled ``convex sets can have interior hot spots'': \url{https://arxiv.org/abs/2412.06344}.
\item The exact supplied literature-check date is \texttt{2026-08-21}. This is supplied dated curation, not an independently verified assessment. It labels the record partially solved and states that log-concavity and the classical fundamental-gap bound are established in major cases, that the hot-spots assertion has counterexamples in sufficiently high dimensions, and that magnetic and Robin extensions remain separate questions. The background expressly calls this dated literature triage derived from the supplied OWR curation and requests independent verification of time-sensitive status and cited claims.
\item The registry labels the statement \texttt{corrected\_verified} and reports that its mathematical content was checked against the official Oberwolfach report, with notation, context, and question form normalized without changing the source-level claim. This verification claim is supplied metadata; it has not been independently checked for this draft.
\end{itemize}
\end{document}
